\section{A stronger certified answer to the natural-density question}
\label{sec:density}

For a set $E\subseteq\N$, write $d(E)$ for its natural density when it exists.
Let $\alpha$ be the supremum of $d(E)$ over sets containing no distinct triple
$a,ar,ar^2$ with rational $r>1$. Let $\beta$ be the analogous supremum when only
integer ratios $r>1$ are forbidden. For instance, $(4,6,9)$ is forbidden by the
first rule and allowed by the second. McNew's published Question 7.4 asks whether
$\beta>\alpha$ \cite{McNew}.

\begin{theorem}[Certified density separation]\label{thm:density-gap}
With these definitions,
\[
 \boxed{\quad \beta-\alpha\ge
 \frac{135102729954660}{10^{18}}>0.0001351.\quad}
\]
\end{theorem}

The strict separation was already asserted and supported by certificates in
the incoming report \cite{priorDensity}. Our contribution is an independent
check of that argument, a stronger finite exchange inequality, a replayable
branch-and-dual certificate, and the sharper constant. This section gives the
complete passage from that finite assertion to natural densities. The result
does not depend on RH, GRH, prime-distribution estimates, or floating-point
optimization. It also does not determine either supremum exactly or settle
McNew's separate conjecture that Rankin's set maximizes $\alpha$.

\subsection{The finite exchange with a boundary charge}
Let $S=\{2^i3^j5^k:i,j,k\ge0\}$ and write $W(A)=\sum_{a\in A}a^{-1}$.
Set
\[
 N=10^6,\qquad \tau=\frac{824341}{10^7},\qquad
 g=\frac{237}{100000},\qquad V=S\cap[1,N].
\]
The set $V$ has 507 elements. Use the 102-element set $B$ in
Table~\ref{tab:density-replacement}, inherited without change from
\cite{priorDensity}.

\begin{lemma}[Exact finite exchange]\label{lem:density-exchange}
The set $B$ contains no integer-ratio geometric triple. Every
rational-ratio-free $A\subseteq S$ satisfies
\begin{equation}\label{eq:density-exchange}
 W(B)-W(A)-\tau W(B\setminus A)\ge g.
\end{equation}
Moreover
\[
 \sum_{r\in\N\setminus S}r^{-2}<\tau.
\]
\end{lemma}

\begin{table}[htbp]\centering\small
\caption{The inherited replacement set $B$. Each listed integer is used once.}
\label{tab:density-replacement}
\begin{tabular}{rrrrrrrrrr}\toprule
1&2&3&5&6&8&10&15&16&24\\
27&40&48&54&60&80&81&90&120&125\\
135&150&162&180&216&250&270&300&375&405\\
432&450&512&625&648&750&900&1000&1024&1080\\
1250&1296&1536&1875&2000&2560&3000&3072&3375&4096\\
5000&5120&6750&7680&8192&10000&10125&12288&13824&15360\\
16875&19683&20480&24576&27000&27648&39366&40960&41472&50625\\
59049&61440&64000&69120&82944&98415&110592&118098&122880&128000\\
138240&157464&192000&196830&207360&221184&295245&314928&320000&331776\\
384000&414720&472392&512000&531441&552960&590490&640000&663552&787320\\
944784&960000&&&&&&&&\\\bottomrule
\end{tabular}
\end{table}

\begin{proof}[Finite verification and why its certificate proves the assertion]
The code enumerates $V$ by multiplication by $2,3,5$ and directly checks
for every $a\in B$ and $2\le r\le\sqrt{N/a}$ that $ar,ar^2$ are not both
in $B$. This exhausts every possible integer-ratio triple in $B$.
Geometric series give $W(S)=15/4$, and exact summation gives
\[
 W(B)=\frac{1928068463873}{680244480000},\qquad
 T:=W(S\setminus V)=\frac{12000753838363}{108839116800000000}.
\]
Give $n\in V$ weight $w_n=(1-\tau\mathbf1_B(n))/n$.
The finite statement to certify is
\begin{equation}\label{eq:density-finite}
 \sum_{n\in I}w_n\le(1-\tau)W(B)-T-g
 \quad\text{for every rational-ratio-free }I\subseteq V.
\end{equation}
Since $W(A\setminus V)\le T$, this rearranges to
\eqref{eq:density-exchange}.

All weights are scaled to integers with
\[
 D=1088391168000000000000000.
\]
In these units the independent-set upper bound is
\[
 H=2827908305897851139120000,
\]
equivalently every deletion set meeting each rational triple must have weight
at least
\[
 Q=999136824871360000000000.
\]
There are 16,446 rational triples in $V$. The first verifier generates them
from endpoints $a<c$ with $ac$ a square and $\sqrt{ac}\in V$. It gives a
finite proof tree with two kinds of valid steps:
\begin{enumerate}
\item At a leaf, assign nonnegative integer masses to remaining triples so
that their total incident mass at each vertex is at most its weight. Every
vertex cover costs at least the total mass: sum the covering inequality for
each triple and then use the vertex capacities. A leaf is certified when this
total reaches the still-required deletion weight.
\item Otherwise choose one uncovered triple. Every cover contains at least
one of its three vertices. Branch on these three possibilities, pay the chosen
vertex's weight, and remove every triple it meets. Proving all three children
proves the parent.
\end{enumerate}
These rules are justified by induction on the finite proof tree. The emitted
JSON records every branch and every leaf's integer dual masses. A separate
replayer checks the capacities, the uncovered-edge conditions, and all three
children at every branch. Generation and replay each visit 496 nodes.

The second verifier uses a different enumeration. It chooses vertices in
increasing order and, whenever a middle term $b$ and an earlier selected term
$a$ are present, excludes the future term $b^2/a$ if integral and in $V$.
An admissible set is represented by one of the include/exclude branches.
The sum of selected and still available weights is an upper bound for any
completion. The search closes every branch at or below $H$ in 416,653 visits.
This verifies \eqref{eq:density-finite} independently of the cover certificate.

Finally, the ratio charge is
\[
 \sum_{r\notin S}r^{-2}=\zeta(2)-\frac{25}{16}.
\]
For $M=10000$, the integral comparison bounds the tail of
$\sum r^{-2}$ by $1/M$. Rounding each of its first $M$ terms upward on
the grid $10^{-18}\Z$ yields the certified upper bound
\[
 \zeta(2)-25/16<
 \frac{82434071848064747}{10^{18}}<\frac{824341}{10^7}.
\]
Every calculation above is an integer or rational operation in
\texttt{code/certify\_density.py}.
\end{proof}

\subsection{Why the exchange survives additional primes}
Let $S_Y$ consist of integers all of whose prime factors are at most $Y$,
where $Y\ge5$. If $E\subseteq S_Y$ is rational-ratio-free, put
$A=E\cap S$. Replace $A$ by $B$. For each newly inserted $b\in B\setminus A$,
delete every $br^2$ belonging to the current set for which the integer $r$
has a prime factor greater than $5$.

The resulting set $F$ is integer-ratio-free. To prove this, consider a newly
created triple. If its ratio is $5$-smooth, the valuations at primes greater
than $5$ are constant along it. Since it contains a newly inserted smooth
integer, it lies entirely in $S$, where $B$ is admissible. If its ratio has
an outside prime, a newly inserted smooth integer can only be its first term;
its final term has been deleted. A triple consisting entirely of retained old
elements was already forbidden in $E$.

The deletion cost is at most
$\tau W(B\setminus A)$, even allowing all positive integer ratios rather
than just $Y$-smooth ratios. Thus Lemma~\ref{lem:density-exchange} gives
\begin{equation}\label{eq:density-uniform}
 W(F)\ge W(E)+g
 \qquad\text{uniformly for every }Y\ge5.
\end{equation}
This uniformity is essential: an improvement on a fixed finite set alone
would have no positive natural-density consequence.

\subsection{From reciprocal weights to natural densities}
Let $U_Q(Y)$ and $U_I(Y)$ be the suprema of $W$ on rational-ratio-free and
integer-ratio-free subsets of $S_Y$. They are finite, since
\[
 W(S_Y)=c_Y^{-1},\qquad
 c_Y=\prod_{p\le Y}(1-p^{-1}).
\]
Equation~\eqref{eq:density-uniform} proves $U_I(Y)\ge U_Q(Y)+g$.

Suppose $E\subseteq\N$ is rational-ratio-free with natural density $d$.
Factor each integer uniquely as $hs$ with $s\in S_Y$ and with $h$ having
no prime factors at most $Y$. Each fixed $h$-fiber of $E$ gives an admissible
smooth set. Therefore
\[
 \sum_{\substack{n\in E\\n\le x}}\frac1n
 \le U_Q(Y)\sum_{\substack{h\le x\\p\mid h\Rightarrow p>Y}}\frac1h.
\]
Partial summation makes the left side $(d+o(1))\log x$. The rough integers
form a periodic coprimality set, so the rightmost sum is
$(c_Y+o(1))\log x$. Taking suprema gives
\begin{equation}\label{eq:density-upper}
 \alpha\le c_Y U_Q(Y).
\end{equation}
No density is assumed to exist within an individual fiber.

For the lower bound choose any \emph{finite} integer-ratio-free set
$F\subset S_Y$. Use all integers $hs$ with $s\in F$ and $h$ squarefree
and rough at $Y$. In any geometric triple the outside prime valuations
would be an arithmetic progression in $\{0,1\}$, hence constant.
The three terms would have the same rough part $h$, contradicting the
admissibility of $F$. The rough squarefree integers have density
\[
 c_Yv_Y,\qquad v_Y=\prod_{p>Y}(1-p^{-2}):
\]
prove this first for finitely many square restrictions by the Chinese
remainder theorem, then bound the omitted integers by
$x\sum_{p>K}p^{-2}$. The finitely many disjoint dilated fibers have
natural density $c_Yv_Y W(F)$.

Finite sets approximate $U_I(Y)$ in reciprocal weight, since the total
smooth mass converges. Taking their supremum consequently gives
$\beta\ge c_Yv_Y U_I(Y)$. Together with \eqref{eq:density-upper},
\begin{align}
 \beta-\alpha
 &\ge c_Y\bigl(U_I(Y)-U_Q(Y)\bigr)
       -c_YU_I(Y)(1-v_Y)\notag\\
 &\ge c_Yg-\sum_{p>Y}p^{-2}.\label{eq:density-bridge}
\end{align}
Here $c_YU_I(Y)\le1$ and
$1-\prod(1-a_p)\le\sum a_p$ for $0\le a_p\le1$.
The construction does not assert a density for an infinite optimizing
smooth set; it only takes suprema after constructing sets whose densities exist.

\subsection{The certified numerical value}
Take $Y=6619$ and sieve through $2,000,000$. Starting with the integer
$10^{18}$, multiply successively by $(p-1)/p$ and round downward for each
$p\le Y$. For $Y<p\le2,000,000$, round $10^{18}/p^2$ upward, and add
$10^{18}/2,000,000$ to dominate the remaining tail even over all integers.
This proves
\[
 c_Y\ge\frac{63728234915742846}{10^{18}},\qquad
 \sum_{p>Y}p^{-2}\le\frac{15933186795650}{10^{18}}.
\]
Substitution of $g=237/100000$ into \eqref{eq:density-bridge}, with one
final downward rounding, proves Theorem~\ref{thm:density-gap}.
The earlier report used $g=23/10000$ and $Y=6841$, obtaining
$130649725461309/10^{18}$. The stronger bound here increases that
certified gap by about $3.4\%$; this is a quantitative improvement to the
existing separation proof, not a new claim of first discovery.

\paragraph{Connection to the source method.}
For integers, $ac=b^2$ is equivalent to a midpoint relation in every
prime-exponent coordinate. The repository's
\src{New/Reduction/Data/Existence.lean}{progression-free coordinate construction}
is one motivation for studying such constraints. Integer ratios impose
nonnegative coordinate increments; rational ratios allow mixed signs.
That asymmetry is precisely what makes the endpoint deletion in the exchange
argument work. The certificate and the density bridge are additional
combinatorial reasoning, not corollaries of the repository's final zero theorem.
